\subsection{庞加莱群的作用与覆叠的态射}

设 B 是连通且局部道路连通的拓扑空间，b 是 B 的一点。

设 E 是 B 的覆叠。由于假设空间 B 连通，若 E 非空，则纤维 \(E_{b}\) 非空（I，第 74 页，命题 4）。根据 III 第 305 页定理 1 的断言 a)，覆叠 E 连通且非空，当且仅当群 \(\pi_{1}(B,b)\) 在 \(E_{b}\) 上传递作用。

命题 2. --- 设 E 和 E' 是 B 的覆叠。

\begin{enumerate}
\def\labelenumi{\alph{enumi})}
\item
设 \(f\colon\mathrm{E}\to\mathrm{E}^{\prime}\) 是 B-态射。映射 \(f_{b}\colon\mathrm{E}_{b}\to\mathrm{E}_{b}^{\prime}\) 由 \(f\) 导出，分别与 \(\pi_{1}(\mathrm{B},b)\) 在 \(\mathrm{E}_{b}\) 和 \(\mathrm{E}_{b}^{\prime}\) 上的作用相容。当且仅当 \(f\) 是同构时，它是双射。
\item
映射 \(f\mapsto f_{b}\) 是从 B-态射集合 \(\mathcal{C}_{\mathrm{B}}(\mathrm{E};\mathrm{E}^{\prime})\)（其元素是从 E 到 \(\mathrm{E}^{\prime}\) 的 B-态射）到 \(\mathcal{F}_{\pi_1(\mathrm{B},b)}(\mathrm{E}_b;\mathrm{E}_b^{\prime})\)（即 \(\pi_1(\mathrm{B},b)\)-态射集合，其元素是从 \(\mathrm{E}_b\) 到 \(\mathrm{E}_b^{\prime}\) 的态射）的双射。
\end{enumerate}

若 \(f\) 是从 E 到 \(\mathrm{E}'\) 的 B-态射，则映射 \(f_b\colon \mathrm{E}_b\to \mathrm{E}_b'\) 是 \(\pi_1(\mathrm{B},b)\)-态射（参见 III，第 305 页）。此外，两个从 E 到 \(\mathrm{E}'\) 的 B-态射在纤维 \(\mathrm{E}_b\) 上相同，则彼此相等（I，第 80 页，命题 6 的推论 3）。

设 \(\varphi\colon\mathrm{E}_{b}\to\mathrm{E}_{b}^{\prime}\) 是 \(\pi_{1}(\mathrm{B},b)\)-集的态射；证明存在一个 B-态射 \(f\)，它从 E 到 \(\mathrm{E}^{\prime}\) 且满足 \(f_{b}=\varphi\)。不妨设空间 E 和 \(\mathrm{E}^{\prime}\) 连通且非空，于是 \(\mathrm{E}_{b}\) 和 \(\mathrm{E}_{b}^{\prime}\) 是齐性 \(\pi_{1}(\mathrm{B},b)\)-集。设 \(x\) 是 \(\mathrm{E}_{b}\) 的一点；它的稳定子 G 是映射 \(\pi_{1}(p,x)\) 的像（III，第 305 页，定理 1）。由于 \(\varphi\) 是 \(\pi_{1}(\mathrm{B},b)\)-态射，群 G 固定点 \(x^{\prime}=\varphi(x)\)，它属于 \(\mathrm{E}_{b}^{\prime}\)，故 G 包含于映射 \(\pi_{1}(p^{\prime},x^{\prime})\) 的像中。由 III，第 308 页的命题 1，存在一个 B-态射 \(f\)，它从 E 到 \(\mathrm{E}^{\prime}\) 且满足 \(f(x)=x^{\prime}\)。映射 \(f_{b}\) 和 \(\varphi\) 是 \(\pi_{1}(\mathrm{B},b)\)-态射，定义在齐性 \(\pi_{1}(\mathrm{B},b)\)-集 \(\mathrm{E}_{b}\) 到 \(\mathrm{E}_{b}^{\prime}\) 之间，且在点 \(x\) 处相同；所以它们相等。

若 \(f\colon \mathrm{E}\to \mathrm{E}'\) 是 B-同构，则映射 \(f_{b}\colon \mathrm{E}_{b}\to \mathrm{E}_{b}^{\prime}\) 是双射。反之，设映射 \(f_{b}\) 是双射，并令 \(g\colon \mathrm{E}^{\prime}\to \mathrm{E}\) 为满足 \(g_{b} = (f_{b})^{-1}\) 的 B-态射。B-态射 \(g\circ f\colon \mathrm{E}\to \mathrm{E}\) 在 \(\mathrm{E}_b\) 上诱导恒等映射；因此由本命题可知 \(g \circ f = Id_{E}\)。同理，\(f \circ g = Id_{E'}\)，这证明 f 是 B-同构。

推论 1. --- 覆叠 E 和 E' 同构，当且仅当 \(\pi_1(\mathrm{B},b)\)-集 \(\mathrm{E}_b\) 和 \(\mathrm{E}_b'\) 同构。

推论 2. --- 设 (E,p) 和 (E',p') 是 B 的连通覆叠。设 \(x\) 是 \(E_{b}\) 的一点，\(x'\) 是 \(E'_{b}\) 的一点。要存在 B-态射 \(g\colon E\to E'\) 使 \(g(x)=x'\)，必要且充分条件是 \(p_{*}(\pi_{1}(E,x))\subset p'_{*}(\pi_{1}(E',x'))\)。这样的态射唯一；并且当且仅当 \(p_{*}(\pi_{1}(E,x))\) 和 \(p'_{*}(\pi_{1}(E',x'))\) 是 \(\pi_{1}(B,b)\) 的相等子群时，它是同构。

由 III，第 305 页定理 1，从 \(\pi_1(B,b)\) 到 \(E_b\) 的映射由 \(\gamma \mapsto x \cdot \gamma\) 定义，经取商诱导出 \(\pi_1(B,b)\)-集 \(p_*(\pi_1(E,x)) \backslash \pi_1(B,b)\) 到 \(E_b\) 的同构。同样，存在唯一的从 \(\pi_1(B,b)\)-集 \(p'_{*}(\pi_1(E',x')) \backslash \pi_1(B,b)\) 到 \(E_b'\) 的同构。由命题 2，存在一个 B-态射 \(g\colon E\to E'\)，使 \(g(x)=x'\)，当且仅当存在一个 \(\pi_1(B,b)\)-集态射，从 \(p_*(\pi_1(E,x)) \backslash \pi_1(B,b)\) 到 \(p'_{*}(\pi_1(E',x')) \backslash \pi_1(B,b)\)，将类 \(p_*(\pi_1(E,x))\) 送到类 \(p'_{*}(\pi_1(E',x'))\)。这样的态射存在，当且仅当 \(p_*(\pi_1(E,x)) \subset p'_{*}(\pi_1(E',x'))\)。由于假设空间 E 连通，它因而是唯一的（I，第 34 页，命题 11 的推论 2）。当且仅当子群 \(p_*(\pi_1(E,x))\) 和 \(p'_{*}(\pi_1(E',x'))\) 是 \(\pi_1(B,b)\) 的相等子群时，它是同构。

推论 3. --- 设 (E,p) 是 B 的连通覆叠，x 是纤维 \(E_{b}\) 的一点。令 N 表示 \(p_{*}(\pi_{1}(E,x))\) 在 \(\pi_{1}(B,b)\) 中的正规化子。对于 N 的每个元素 \(\gamma\)，存在唯一的 B-自同构 g，使得 \(g(x) = x \cdot \gamma\)；并且映射 \(\gamma \mapsto g\) 经取商定义了从 \(\mathrm{N}/p_{*}(\pi_{1}(\mathrm{E},x))\) 到 \(\mathrm{Aut}_{\mathrm{B}}(\mathrm{E})\) 的群同构。

设 \(\gamma \in \pi_1(\mathrm{B}, b)\)；有 \(p_*(\pi_1(\mathrm{E}, x)) = \operatorname{Int}(\gamma)(p_*(\pi_1(\mathrm{E}, x \cdot \gamma)))\)（III，第 305 页，定理 1，c)）。由推论 2，要存在 B-自同构 \(g\) 使 \(g(x) = x \cdot \gamma\)，必要且充分条件是 \(\gamma\) 属于 N。这样的同构因而唯一。记 \(\alpha: \mathrm{N} \to \mathrm{Aut}_{\mathrm{B}}(\mathrm{E})\) 为由此定义的映射 \(\gamma \mapsto g\)。

设 \(\gamma\) 和 \(\gamma'\) 是 N 的两个元素。置 \(g = \alpha(\gamma)\)、\(g' = \alpha(\gamma')\)；根据关系式 (4)（III，第 305 页）和 (1)（第 304 页），有 \(g(g'(x)) = g(x \cdot \gamma') = g(x) \cdot \gamma' = (x \cdot \gamma) \cdot \gamma' = x \cdot (\gamma\gamma')\)。因此，\(\alpha\) 是群同态。要有 \(\alpha(\gamma) = \mathrm{Id}_{\mathrm{E}}\)，必要且充分条件是 \(x \cdot \gamma = x\)，因为 \(\alpha(\gamma)\) 唯一；也就是说，\(\gamma \in p_*(\pi_1(E, x))\)（III，第 305 页，定理 1，b)）。最后，若 \(g\) 是 E 的 B-自同构，则存在一条道路 \(c\)，它在 E 中连接 \(x\) 与 \(g(x)\)（E 是道路连通的），于是 \(g(x) = x \cdot \gamma\)，其中 \(\gamma\) 是道路 \(p \circ c\) 的类。这证明同态 \(\alpha\) 是满射。

群 \(\mathrm{Aut}_{\mathrm{B}}(\mathrm{E})\) 在纤维 \(\mathrm{E}_b\) 上作用，并被识别为齐性 \(\pi_1(\mathrm{B},b)\)-集 \(\mathrm{E}_b\)（右作用）的自同构群（参见 A，I，第 56 页，命题 5 和 6）。

推论 4. --- 设 (E, p) 是 B 的连通覆叠，x 是纤维 \(E_{b}\) 的一点。要使 E 成为 B 的伽罗瓦覆叠，必要且充分条件是 \(p_{*}(\pi_{1}(E, x))\) 是 \(\pi_{1}(B, b)\) 的正规子群。此时群 \(\mathrm{Aut}_{\mathrm{B}}(\mathrm{E})\) 同构于商群 \(\pi_{1}(\mathrm{B}, b)/p_{*}(\pi_{1}(\mathrm{E}, x))\)。

若 \(p_{*}(\pi_{1}(\mathrm{E},x))\) 是 \(\pi_1(\mathrm{B},b)\) 的正规子群，则根据推论 3，群 \(\mathrm{Aut}_{\mathrm{B}}(\mathrm{E})\) 在纤维 \(\mathrm{E}_b\) 上传递作用，所以 E 是 B 的伽罗瓦覆叠（I，第 102 页，定理 2）。逆命题是 III，第 305 页定理 1 的断言 d)。最后一个断言由推论 3 得出。

